State the essential results of Elimination theory (affine case)


In the succeeding the lex order is always applied.

Given \( I = \langle f_1, \dots, f_s \rangle \subseteq k[x_1, \dots, x_n] \) the \( l \)-th elimination ideal
refers to the ideal \( I_l \) defined by \( I_l = I \cap k[x_{l+1}, \dots, x_n] \).

Thus \( I_l \) consists of all consequences of \( f_1 = \dots = f_s = 0 \) where the variables \( x_1, \dots, x_l \) have been eliminated.


Elimination Theorem (Groebner basis form)
Let \( I  \subseteq k[x_1, \dots, x_n] \) be an ideal with Groebner basis \( G \)
then \( G_l = G \cap k[x_{l+1}, \dots, x_n] \) is the Groebner basis of \( I_l \).

The reverse direction of extending partial solutions has some pitfalls one needs to circumvent.
Essentially after finding all solutions \( \{(a_{l+1}, \dots, a_n) \in V(I_l) \} \), we determine \( I_{l-1} = \{g_1, \dots, g_r\} \) and try to find
all \( x_l \) with \( \{(x_l, a_{l+1}, \dots, a_n) \in V(I_{l-1}) \} \), which might seem trivial on first sight. 
However, consider \( xy - 1 = 0, xz -1 = 0 \), then \( I_1 = \langle y - z \rangle \) and \( V(I_1) = \{ (y, z) | y = z = \text{ const } c \in k\} \), 
extending this to \( V(I_0) = V(I) \) is valid for all \( (\frac{1}{c}, c, c) \) except for \( c = 0 \).

INSERT IMAGE FROM PAGE 130 HERE


Extension theorem
Let \( I = \langle f_1, \dots, f_s \rangle \subset k[x_1, \dots, x_n] \), \( k = \bar{k} \)
Write each \( f_i \in I_1 \) as \( f_i = h_i(x_2, \dots, x_n) x_1^{N_i} + \text{ terms in which } x_1 \text{ has degree } < N_i \),
where \( N_i \geq 0 \) and \( h_i \in \mathbb{C}[x_2, \dots, x_n] \) is nonzero. Suppose that we have
a partial solution \( (a_2, \dots, a_n) \in V(I_1) \). If \( (a_2, \dots, a_n) \nin V(c_1, \dots, c_s) \),
there there exists \( \alpha_1 \in \mathbb{C} \) s.t. \( (a_1, \dots, a_n) \in V(I) \).


This holds analogously for \( I_l \) and \( I_{l-1} \).
For non algebraically closed fields this works as long as one checks if extended partial solutions would stay in \( k \).
Note that the minimal extension property of this theorem is the content of the closure theorem.


State the closure theorem and the geometric version of the extension theorem

Let \( \pi_l : (a_1, \dots, a_n) \mapsto (a_{l+1}, \dots, a_n) \) denote the projection

Geometric extension theorem
Given \( F = V(f_1, \dots, f_s) \subseteq k^n \), \( k \) alg. closed,
let \( h_i \in k[x_1, \dots, x_n] \) as in the previous extension theorem.
We have the equality
\( V(I_1) = \pi_1(F) \cup (V(h_1, \dots, h_s) \cap V(I_1)) \)
In particular if there's some \( i \) s.t. \( h_i \in k \) (so without zeros) then \( \pi_1(V) = V(I_1) \).

This tells us that \( \pi_1(F) \) fills up the affinte variety \( V(I_1) \) except for a part that lies 
in \( V(c_1, \dots, c_s) \). It is not easy to determine how large \( V(c_1, \dots, c_s) \) is.

Closure theorem
Given \( F = V(f_1, \dots, f_s) \subseteq k^n \), \( k \) alg. closed, then
- \( V(I_l) \) is the smallest affine variety containing \( \pi_l(V) \subseteq k^{n-l} \)
- When \( F \neq 0 \) there is an affine variety \( W \subsetneq V(I_l) \) s.t. \( V(I_l) \setminus W \subseteq \pi_l(F) \).
  Put differently \( V(I_l) \) is the Zariski closure of \( \pi_l(F) \).




State the theorems for implicitization as consequences of Elimination theory (affine case)

Recall that polynomial parametrization \( H : k^m \to k^n \) tries to express a variety in \( (x_1, \dots, x_n) \) via 
a set of parameters \( (t_1, \dots, t_m) \) as \( x_i = h_i(t_1, \dots, t_m) \)

Theorem (Polynomial Implicitization). 
If \( k \) is an infinite field, let \( H : k^m \to k^n \) be the function determined by the polynomial parametrization. 
Let \( I \) be the ideal \( I = \langle x_1 - h_1, \ldots, x_n - h_n \rangle \subseteq k[t_1, \ldots, t_m, x_1, \ldots, x_n] \) and 
let \( I_m = I \cap k[x_1, \ldots, x_n] \) be the \( m \)-th elimination ideal. 
Then \( V(I_m) \) is the smallest variety in \( k^n \) containing \( H(k^m) \).

Note that this makes use of the graph of \( H \) as variety.
The theorem gives the following implicitization algorithm for polynomial parametrizations: 
if we are given \( x_i = h_i(t_1, \dots, t_m) \) for polynomials \( h_i \in k[t_1, \dots, t_m] \), 
consider the ideal \( I = \langle x_1 - h_1, \dots, x_n - h_n \rangle \) and compute a Groebner
basis with respect to a lexicographic ordering where every \( t_i > x_j \) for all \( i,j \).
By the Elimination Theorem, the elements of the Groebner basis not involving \( t_i \)
form a basis of \( I_m \), and by the Theorem, they define the smallest variety in \( k^n \) containing the parametrization.

Now allow for the \( h_i = \frac{f_i}{g_i} \) to be rational.

Theorem (Rational Implicitization).
If \( k \) is an infinite field, let \( H : k^m \setminus W \to k^n \) be the function determined by the rational parametrization. 
Let \( J \) be the ideal 
\( J = \langle g_1 x_1 - f_1, \ldots, g_n x_n - f_n, 1 - g y \rangle \subseteq k[y, f_1, \ldots, f_n, t_1, \ldots, t_m, x_1, \ldots, x_n] \), 
where \( g = g_1 \cdot g_2 \cdots g_n \) and \( W = V(g) \). 
Also let \( J_{1+m} = J \cap k[x_1, \ldots, x_n] \) be the \( (1+m) \)-th elimination ideal. 
Then \( V(J_{1+m}) \) is the smallest variety in \( k^n \) containing \( F(k^m \setminus W) \).

The Theorem implies the following implicitization algorithm for rational parametrizations: 
if we have \(x_i = f_i/g_i\) for polynomials \(f_1, g_1, \ldots, f_n, g_n \in k[t_1, \ldots, t_m]\), 
consider the new variable \(y\) and set \(J = \langle g_1 x_1 - f_1, \ldots, g_n x_n - f_n, 1 - g y \rangle\), 
where \(g = g_1 \cdots g_n\). Compute a Groebner basis with respect to a lexicographic ordering where 
\(y\) and every \(t_i\) are greater than every \(x_i\). 
Then the elements of the Groebner basis not involving \(y, t_1, \ldots, t_m\) define the smallest variety in \(k^n\) containing the parametrization. 
This algorithm is due to Kalkbrenner.



State Hilbert's Nullstellensaetze and some important Corollaries


Weak Nullstellensatz
Let \( k \) be an algebraically closed field and \( I \subseteq k[x_1, \dots, x_n] \) be an ideal satisfying \( V(I) = \emptyset \). 
Then \( I = k[x_{1}, \dots, x_n] \).
(Hard to prove)

Strong Nullstellensatz
Let \( k \) be an algebraically closed field.
If \( f, f_1, \dots, f_s \in k[x_1, \dots, x_n] \) then \( f \in I(V(f_1, \dots, f_s)) \) if and only if
\( f^m \in \langle f_1, \dots, f_s \rangle \) for some integer \( m \geq 1 \).

Ansatz
By assumption
\[
    f^m = \sum_{i=1}^s h_i f_i
\]
for some \( h_i \in k[x_1, \dots, x_n] \), consider the ideal \( \tilde{I} = \langle f_1, \dots, f_s,  1 - yf \rangle  \)
in the ring \( k[x_1, \dots, x_n, y] \). By case consideration one realizes that \( V(\tilde{I}) = \emptyset \).
By the weak Nullstellensatz we have some \( h_i \in k[x_1, \dots, x_n] \) with
\( 1 = h_{s+1}(1 - yf) \sum_{i=1}^s h_i f_i  \), now set \( y = \frac{1}{f} \) then rewriting yields
\( 1 = \sum_{i=1}^s h_i(x_1, \dots, \frac{1}{f}) f_i(x_1, \dots, \frac{1}{f}) \) 
thus multiplying by some power of \( f^m \) cancels all occurences of \( \frac{1}{f} \) on the RHS but leaves some power of \( f \) on the LHS.

A more idiomatic way of expressing this result is
Let \( k \) be an algebraically closed field. If \( I \subseteq k[x_1, \dots, x_n] \) is an ideal then 
\( I(V(I)) = \sqrt{I} \).
Ansatz
\( \sqrt{I} \subseteq I(V(I)) \) follows from \( f^m \) vanishing on \( V(I) \) implies that \( f \) vanishes on \( V(I) \).
Conversely, take \( f \in I(V(I)) \). Then, by definition, f vanishes on \( V(I) \). By Hilbert’s Nullstellensatz,
there exists an integer \( m \geq \sqrt{I} \) such that \( f^m \in I \), so \( f \in \sqrt{I} \).

Corollary 12. If \( k \) is an algebraically closed field, then there is a one-to-one correspondence 
between points of \( k^n \) and maximal ideals of \( k[x_1, \dots, x_n] \).

Theorem 11. If \( k \) is an algebraically closed field, then every maximal ideal of
\( k[x_1, \dots, x_n] \) is of the form \( \langle x_1 - a_1, \dots, x_n - a_n \rangle \) for some \( a_1, \dots, a_n \in k \).

Proposition 13. A variety \( V \) is irreducible if and only if for every variety \( W \subseteq V \),
the difference \( V \setminus W \) is Zariski dense in \( V \).


State the Ideal-Variety correspondence


Let \( k \) be an arbitrary field.

(i) The maps
\[
\text{affine varieties} \xrightarrow{\mathrm{I}} \text{ideals}
\]
and
\[
\text{ideals} \xrightarrow{\mathrm{V}} \text{affine varieties}
\]
are inclusion-reversing, i.e., if \( I_1 \subseteq I_2 \) are ideals, then \( \mathrm{V}(I_1) \supseteq \mathrm{V}(I_2) \) and, similarly, if \( V_1 \subseteq V_2 \) are varieties, then \( \mathrm{I}(V_1) \supseteq \mathrm{I}(V_2) \).

(ii) For any variety \( V \),
\[
\mathrm{V}(\mathrm{I}(V)) = V,
\]
so that \( \mathrm{I} \) is always one-to-one. On the other hand, any ideal \( I \) satisfies
\[
\mathrm{V}(\sqrt{I}) = \mathrm{V}(I).
\]

(iii) If \( k \) is algebraically closed, and if we restrict to radical ideals, then the maps
\[
\text{affine varieties} \xrightarrow{\mathrm{I}} \text{radical ideals}
\]
and
\[
\text{radical ideals} \xrightarrow{\mathrm{V}} \text{affine varieties}
\]
are inclusion-reversing bijections which are inverses of each other.




Define the Zariski topology and related notations

Consider a subset \( E \subset A \) of a commutative unital ring \( A \) 
and the associated ideal \( \mathfrak{a} = (E) \subset A \) generated by \( E \) in \( A \). 
Define
\[
\begin{aligned}
V(E) &= \{ x \in \mathrm{Spec}\, A ;\, f(x) = 0 \text{ for all } f \in E \} \\
     &= \{ x \in \mathrm{Spec}\, A ;\, E \subset \mathfrak{p}_x \}
\end{aligned}
\]
it is also referred to as the variety of \( E \), which explains the usage of the letter \( V \). 
For \( f \in A \) we will apply the notations
\[
V(f) = \{ x \in \mathrm{Spec}\, A ;\, f(x) = 0 \},
\]
\[
D(f) = \mathrm{Spec}\, A - V(f) = \{ x \in \mathrm{Spec}\, A ;\, f(x) \ne 0 \}
\]
\( D(f) \) is also known as the domain of \( f \), which explains the usage of the letter \( D \). 

Proposition:
Consider subsets \( E, E' \), and a family of subsets \( (E_\lambda)_{\lambda \in \Lambda} \) of a ring \( A \). Then:
(i) \( V(0) = \mathrm{Spec}\, A \), \( V(1) = \emptyset \).
(ii) \( E \subset E' \implies V(E) \supset V(E') \).
(iii) \( V \left( \bigcup_{\lambda \in \Lambda} E_\lambda \right) = V \left( \sum_{\lambda \in \Lambda} E_\lambda \right) = \bigcap_{\lambda \in \Lambda} V(E_\lambda) \).
(iv) \( V(E E') = V(E) \cup V(E') \), where \( E E' = \{ f f' ;\, f \in E,\, f' \in E' \} \).
(v) \( V(E) = V(\sqrt{(E)}) \).

Let \( A \) be a commutative unital ring and \( X = \text{Spec } A \) its spectrum. There exists a
unique topology on \( X \), the so-called Zariski topology, whose closed sets consist
precisely of the subsets of type \( V(E) \subset X \) for some subset \( E \subset A \).
Furthermore:
(i) The sets of type \( D(f) \) for \( f, f' \in A \) are open and satisfy the intersection
relation \( D(f) \cap D(f') = D(f f') \).
(ii) Every open subset of \( X \) is a union of sets of type \( D(f) \). In particular,
the latter form a basis of the Zariski topology on \( X \) (which is the reason why
the sets of type \( D(f) \) are referred to as the basic open subsets of X).
The Zariski topology is Hausdorff and for any \( f \in A \), \( D(f) \) is quasi-compact.

In particular \( V(E) = \bigcap_{f \in E} V(f) \) and any open \( U \subset X \) can be expressed as \( U = \bigcup_{f \in E} D(f) \).
Note that \( D(E) = \bigcap_{f \in E} D(f) \) which is again open if \( (E) \) has a finite basis.

The somewhat inverse notion to \( V \) is defined as 
\( I(Y) = \{ f \in A ; f(y) = 0 for all y \in Y \} = \bigcap_{y \in Y} \mathfrak{p}_y \).

Let A be a ring.
(i) Let \( Y \subset Y' \subset \text{Spec } A \) be subsets. Then \( I(Y) \supsect I(Y') \).
(ii) \( I(\bigcup_{\lambda \in \Lambda} Y_\lambda) = \bigcap_{\lambda \in \Lambda} \) for any family \( \bigcup_{\lambda \in \Lambda} Y_\lambda) \)
    of subsets of \( \text{Spec } A \).

Hilbert's Nullstellensatz then expresses \( I(V(I)) = \sqrt{I} \) and for any \( Y \subset \text{Spec } A \), \( V(I(Y)) \) coincides
with the Zariski closure of \( Y \). Thus for closed \( Y \), \( V \) and \( I \) are inverse and inclusion-reversing.
In particular
\( I(\bigcup_{\lambda \in \Lambda } = \bigcap_{\lambda \in \Lambda} I(Y_\lambda) \) for any \( Y \subset \text{Spec } A  \)
and for closed \( Y_\lambda \) it holds
\( I(\bigcap_{\lambda \in \Lambda } = \sqrt{ \sum_{\lambda \in \Lambda} I(Y_\lambda)} \).

For any ideal \( \mathfrak{a} \subset A \) it holds that the homomorphism \( \pi : A \to A / \mathfrak{a} \)
introduces a map 
\( \tilde{\pi} : \text{Spec} A / \mathfrak{a} \to \text{Spec} A \) via \( \mathfrak{p} \mapsto \pi^{-1}(\mathfrak{p}) \)
which is a homomorphism if restricting the codomain to \( V(\mathfrak{a}) \), where \( V(\mathfrak{a}) \) has 
the subspace topology of \( \text{Spec } A \).



Define and give algorithms for arithmetic operations on varieties

\( I + J \) is the smallest ideal containing \( I, J \) and for \( I = \langle f_1, \dots, f_i \rangle, J = \langle g_1, \dots, g_j \rangle \)
\( I + J = \langle f_1, \dots, f_i, g_1, \dots, g_j \rangle \).
It follows that \( V(I + J) = V(I) \cap V(J) \).

Similarly for \( IJ \) we get \( IJ = \langle f_1 g_1, f_1 g_2, \dots, f_1, g_j, \dots, f_i g_j \rangle \) on the generators
and \( V(IJ) = V(I) \cup V(J) \).

To compute intersections of ideals (which are ideals themselves) we can make use of the fact 
\( I \cap J = (tI + (1 - t)J) \cap k[x_1, \dots, x_n] \).

The above result and the Elimination Theorem lead to the following algorithm for computing intersections of ideals: 
if \( I = \langle f_1, \ldots, f_r \rangle \) and \( J = \langle g_1, \ldots, g_s \rangle \) are ideals in \( k[x_1, \ldots, x_n] \), 
we consider the ideal
\[
\langle t f_1, \ldots, t f_r,\, (1-t) g_1, \ldots, (1-t) g_s \rangle \subseteq k[x_1, \ldots, x_n, t]
\]
and compute a Gröbner basis with respect to lex order in which \( t \) is greater than the \( x_i \). 
The elements of this basis which do not contain the variable \( t \) will form a basis (in fact, a Gröbner basis) of \( I \cap J \).

Define the least common multiple as 
\( \text{lcm}(f, g) = h_1^{\max(a_1, b_1)} \cdots h_l^{\max(a_l, b_l)} f_{l+1}^{a_{l+1}} \cdots g_{l+1}^{b_{l+1}} \),
where \( f_i, g_j \) are the irreducible factors of \( f, g \) appearing with power \( a_i, b_j \) respectively and \( h_i \) those which are
common to both of them.

For principal ideals \( I = \langle f \rangle, J = \langle g \rangle \),
the intersection can then be convienently be expressed as \( I \cap J = \text{lcm}(f,g) \).

Also the relation \( \text{lcm}(f, g) \cdot \text{gcd}(f, g) = fg \) allows for an explicit way of determining the \( gcd \) of \( 2 \) polynomials.

Last geometrically the intersection satisfies \( V(I \cap J) = V(I) \cup V(J) \).


Define the ideal quotient and its algebro-geometric relevance

If \( I, J \subset R \) are ideals in a ring \( R \) then the ideal quotient is the ideal
\( I : J = \{ f \in R | fg \in I, \forall  g \in I \} \).

Proposition
(i) If \( V, W \) are varieties in \( k^n \) then
\( V = (V \cap W) \cup (\overline{V \setminus W}) \)
(ii) If \( I, J \subset k[x_1, \dots, x_n] \) are ideals then
\( V(I) = V(I + J) \cup V(I : J) \)
and \( \overline{V(I) \setminus V(J)} \subseteq V(I:J) \)


The saturation of \( I \) with respect to \( J \) is the set
\( I : J^\infty = \{ f \in R | \forall g \in J, \exist N \geq 0, fg^N \in I \} \).
It holds \( I : J^\infty = \bigcup_{n \in \mathbb{N}} I : J^n \).

For \( R = k[x_1, \dots, x_n] \) it holds
(i) \(I \subseteq I : J \subseteq I : J^\infty \)
(ii) \( I : J^\infty = I : J^N \) for some sufficiently large \( N \).
(iii) \( \sqrt{I : J^\infty} = \sqrt{I} : J \).

Moreover 

(i) If \( V, W \) are varieties in \( k^n \) then
\( V = V(I+J) \cup V(I : J^\infty) \)
(ii) 
\( \overline{V(I) \setminus V(J)} \subseteq V(I: J^\infty) \)
and if \( k \) is alg. closed this is an equality.

For multiple ideals we have

\( I : \sum_{i=1}^r J_i = \bigcap_{i=1}^r (I : J_i) \)
\( I : \left(\sum_{i=1}^r J_i \right)^\infty = \bigcap_{i=1}^r (I : J_{i}^\infty) \)


To actually compute it one can use the following result

Theorem. Let \( I \) be an ideal and \( g \) an element of \( k[x_1, \ldots, x_n] \). Then:

(i) If \(\{h_1, \ldots, h_p\}\) is a basis of the ideal \( I \cap \langle g \rangle \), then \(\{h_1 / g, \ldots, h_p / g\}\) is a basis of \( I : g \).
(ii) If \(\{f_1, \ldots, f_s\}\) is a basis of \( I \) and \( \tilde{I} = \langle f_1, \ldots, f_s, 1 - yg \rangle \subseteq k[x_1, \ldots, x_n, y] \), where \( y \) is a new variable, then
\[
I : g^{\infty} = \tilde{I} \cap k[x_1, \ldots, x_n].
\]
Furthermore, if \( G \) is a lex Gröbner basis of \( \tilde{I} \) for \( y > x_1 > \cdots > x_n \), then \( G \cap k[x_1, \ldots, x_n] \) is a basis of \( I : g^{\infty} \).

algorithm for computing a basis of an ideal quotient: 
Given \( I = \langle f_1, \ldots, f_r \rangle \) and \( J = \langle g_1, \ldots, g_s \rangle \), to compute a basis of \( I : J \), 
we first compute a basis for \( I : g_i \) for each \( i \). 
In view of the previous theorem, this means computing a basis \( \{h_1, \ldots, h_p\} \) of 
\( \langle f_1, \ldots, f_r \rangle \cap \langle g_i \rangle \). 
Recall that we do this via the algorithm for computing intersections of ideals. 
Using the division algorithm, we divide each of basis element \( h_j \) by \( g_i \) to get a basis for \( I : g_i \) by part (i) of the Theorem. 
Finally, we compute a basis for \( I : J \) by applying the intersection algorithm \( s-1 \) times, computing first a basis for 
\( I : \langle g_1, g_2 \rangle = (I : g_1) \cap (I : g_2) \), 
then a basis for \( I : \langle g_1, g_2, g_3 \rangle = (I : \langle g_1, g_2 \rangle) \cap (I : g_3) \), and so on.

Similarly, we have an algorithm for computing a basis of a saturation: 
Given \( I = \langle f_1, \ldots, f_r \rangle \) and \( J = \langle g_1, \ldots, g_s \rangle \), 
to compute a basis of \( I : J^\infty \), we first compute a basis for \( I : g_i^\infty \) for each \( i \) 
using the method described in part (ii) of the Theorem. Then by the summation property of quotient ideal, 
we need to intersect the ideals \( I : g_i^\infty \), which we do as above by applying the intersection algorithm \( s-1 \) times.

